%%%%%%%%%%%%%%%%%%%%%%%%%%%%%%%%%%%%%%%%%%%%%%%%%%%%%%%%%%%%
\section{Global Amplification Barrier}
\label{sec:global_amplification_barrier}
%%%%%%%%%%%%%%%%%%%%%%%%%%%%%%%%%%%%%%%%%%%%%%%%%%%%%%%%%%%%

We formalize a structural limitation of locality-bounded refinement systems:
local refinement cannot synthesize global invariants without paying linear
EntropyDepth cost.

\subsection{Global Invariants}

\begin{definition}[Global Invariant]
Let $\mathcal{G}_n$ be a family of graphs with $|V|=n$.
A function
\[
\Phi : \mathcal{G}_n \to \{0,1\}
\]
is called \emph{global} if for every fixed radius $r$
there exist graphs $G,H \in \mathcal{G}_n$ such that:

\begin{enumerate}
\item For every vertex $x\in G$ there exists $y\in H$ with
\[
\Ball_r^G(x) \cong \Ball_r^H(y),
\]
and vice versa (local indistinguishability),

\item $\Phi(G) \ne \Phi(H)$.
\end{enumerate}
\end{definition}

Thus $\Phi$ is not determined by bounded-radius neighborhoods.

Examples include connectivity, bipartiteness of large girth expanders,
parity of certain cycle counts, and global expansion properties.

\subsection{Chronos Refinement Setting}

Assume the refinement–coarsening process satisfies
FLC$(\Delta,r,k)$ with $\Delta,r$ fixed independently of $n$.

Let $H_t = \log |D_t|$ denote transcript entropy.

Define EntropyDepth:
\[
\ED(G) := \min \{ t : |D_t| = 1 \}.
\]

\subsection{Amplification Principle}

Intuitively, extracting $\Phi$ requires collapsing
all graphs in the same $\Phi$-class to a single transcript value,
while separating distinct $\Phi$-classes. Because local refinement operations
can only propagate information across a distance proportional to the refinement
depth, a global invariant can be resolved only after the execution trace has
reached the scale at which the separating information becomes locally visible.

\subsection{The Barrier Theorem}

We isolate the exact hypothesis needed for a linear lower bound.

\begin{definition}[Linear Local-Indistinguishability Witness]
A global invariant $\Phi$ admits a \emph{linear local-indistinguishability
witness} if there exists a constant $c>0$ and infinitely many $n$ such that
for each such $n$ there are graphs $G_n,H_n\in\mathcal{G}_n$ satisfying
\[
\Phi(G_n)\neq \Phi(H_n),
\]
and for every $t<c n$, the graphs $G_n$ and $H_n$ are indistinguishable by
the refinement transcript up to depth $t$.
\end{definition}

\begin{theorem}[Conditional Global Amplification Barrier]
Assume the refinement process satisfies $\mathrm{FLC}(\Delta,r,k)$ with
constant locality parameters. If $\Phi$ admits a linear
local-indistinguishability witness, then there exists a constant $c>0$ such
that every refinement sequence correctly computing $\Phi$ satisfies
\[
\ED(G_n)\ge c n
\]
for infinitely many $n$.
\end{theorem}

\begin{proof}
Let $G_n,H_n$ be a linear local-indistinguishability witness for $\Phi$.
Suppose that a correct refinement sequence has $\ED(G_n)<c n$.

By $\mathrm{FLC}(\Delta,r,k)$, the transcript at depth $t$ depends only on
bounded-radius data propagated through depth $t$. Hence, for every $t<c n$,
the refinement transcripts of $G_n$ and $H_n$ agree.

Taking $t=\ED(G_n)$, the transcript has already collapsed to a single value
on $G_n$. Since the transcript of $H_n$ is identical at the same depth, the
same collapsed value is assigned to $H_n$. Therefore the refinement sequence
cannot distinguish $G_n$ from $H_n$, contradicting
\[
\Phi(G_n)\neq \Phi(H_n).
\]
Thus $\ED(G_n)\ge c n$ for infinitely many $n$.
\end{proof}
